\phantomsection\label{app:ai}
\subsection*{B.1 使用范围与边界}

本文在以下环节使用了人工智能工具（下称“该工具”）辅助：附件数据的核对与统计、文献检索与思路整理、
公式与能耗关系的复核、程序编写与调试、以及初稿文字的整理与润色。研究主线与建模思路的确定、
模型的建立与取舍、关键假设的设定、结果分析与结论，均由参赛队独立完成；该工具不替代参赛队的
独立判断与核心创新。正文呈现的数学表述、公式与文字均为参赛队自己的语言；凡引自文献或公开资料
之处，均在正文引用处与参考文献中列明。

\subsection*{B.2 工具信息}

\begin{table}[H]\centering
\begin{tabular}{ll}\toprule
名称 & OpenAI Codex \\ \midrule
版本/型号 & GPT-5 \\ \midrule
开发机构/公司 & OpenAI \\ \midrule
版本发布日期 & 2025 年 8 月 7 日 \\ \midrule
使用方式 & 通过对话式接口辅助程序编写、数值核对与文本整理 \\ \midrule
对应文献 & 参考文献\cite{codexai} \\ \bottomrule
\end{tabular}
\caption{本文所用人工智能工具的信息}
\label{tab:ai_tool}
\end{table}

\subsection*{B.3 输入内容与后续处理策略}

该工具的输入以题目附件、题面规则与本队的建模设想为主，典型输入包括“按题给的等效航程与返航安全
余量给出各机型安全载荷的计算程序”“核对某公式的量纲与单位一致性”“把某段说明压缩到指定字数”等。
该工具的输出一律不作为最终结论直接采用，均按以下策略处理：

\begin{enumerate}
  \item 逐条比对该输出与题面参数、单位及题设规则，剔除与题目不一致的部分；
  \item 由参赛队独立推导或复核其中的模型与公式，凡无法给出出处的一律不采用；
  \item 用随附程序重新计算全部数值，并与论文中的数值逐一回代比对；
  \item 正文文字由参赛队重新组织与改写，不直接采用工具生成的文字作为最终表述；
  \item 全部程序与结果纳入版本管理，保留修改记录以备核查。
\end{enumerate}

程序实现所依赖的开发框架与开源软件为：Python 3.13，以及 NumPy、SciPy、openpyxl、rasterio、
pyproj、Affine、OR-Tools（CP-SAT）与 Matplotlib（版本见 \path{requirements.lock}）；论文由
XeLaTeX 与 ctex 宏包排版，示意图由 TikZ 绘制。算法逻辑、假设条件、参数与超参数分别见第 2 至
第 9 章相应小节，此处不再重复。

\subsection*{B.4 合规说明}

\begin{enumerate}
  \item 数据分析结果处已按要求添加说明，见第 2 章开头的“人工智能辅助说明”以及第 5、6 章的结果小节；
  \item 程序实现的工具使用情况统一在本附录 B.3 中说明；\path{code/} 目录下的程序与
        \path{requirements.lock} 记录了全部实现细节，可供核查；
  \item 论文中的数学模型与公式均有文献来源或由参赛队推导，引用见参考文献；本文未使用来源无法确认的
        人工智能生成模型或公式作为结论依据；
  \item 正文文字由参赛队重新组织与改写，未直接采用工具生成的文字作为最终表述。
\end{enumerate}
